An edge $e$ of a Fano polygon $P$ determines the cone $\sigma$ over $e$, and hence a torus-fixed point in the toric variety $X_P$. A neighbourhood of this point is the affine toric variety $X_\sigma$ determined by $\sigma$, that is, the cyclic quotient singularity $\frac{1}{b}(1,c-1)$ where the rays $\{u,v\}$ of $\sigma$ are such that $u$,~$v$, and $\frac{c-1}{b} u + \frac{1}{b}v$ generate $N$. Write $b = dr$ and $c = dk$ where $d = \gcd\{b,c\}$. Then $X_\sigma$ is the cyclic quotient singularity $\frac{1}{dr}(1,kd-1)$. Furthermore, $r$ is the lattice height of $e$ above the origin, and $d$ is the lattice length of $e$.

When $d=nr$ for some $n\in\mathbb{Z}$, the singularity $X_\sigma$ is called a \emph{T-singularity} and $\sigma$ is called a \emph{$T$-cone}. If $n=1$ then $X_\sigma$ is a \emph{primitive} $T$-singularity and $\sigma$ is a \emph{primitive} $T$-cone. $T$-singularities are qG-smoothable~\textup{[28]}. A general $T$-singularity $\frac{1}{n r^2}(1, knr-1)$ admits a crepant partial resolution into $n$ primitive $T$-singularities $\frac{1}{r^2}(1, kr-1)$; this amounts to the fact that, since $d = n r$, the cone $\sigma$ can be decomposed as a union of $n$ primitive $T$-cones -- see figure 2.
When $d = s$ with $0 < s < r$, the singularity $X_\sigma$ is called an \emph{R-singularity} and $\sigma$ is called an $R$-cone. $R$-singularities are qG-rigid. When $d = nr+s$ for some non-negative $n \in \mathbb{Z}$ and $0 < s < r$, the singularity $X_\sigma = \frac{1}{dr}(1, kd-1)$ admits a (non-unique) crepant partial resolution into $n$ primitive $T$-singularities and the single $R$-singularity $\frac{1}{sr}(1,ks-1)$; this corresponds to the (non-unique) decomposition of $\sigma$ into $n$ primitive $T$-cones and one $R$-cone -- see figure 3.
\noindent The singularity $X_\sigma$ qG-deforms to the R-singularity $\frac{1}{sr}(1,ks-1)$. We define the \emph{residue} of~$\sigma$ to be
\[
\operatorname{res}(\sigma)=\begin{cases}
\emptyset&\text{if $s=0$}\\
\frac{1}{sr}(1,ks-1)&\text{otherwise,}
\end{cases}
\]
and the \emph{singularity content} of~$\sigma$ to be the pair $(n,\operatorname{res}(\sigma))$.


\par\noindent\textbf{Definition 3.2.} 
Let $P$ be a Fano polygon with edges $e_1,\ldots,e_m$, and suppose that the singularity content of the cone over $e_i$ is $(n_i, S_i)$. Let $n = n_1 + \cdots + n_m$, and let $\mathcal{B}$ be the multiset containing those $S_i$ such that $S_i \ne \emptyset$, counted with multiplicity. The \emph{singularity content} of $P$ is
\[
\operatorname{SC}(P)=(n, \mathcal{B}).
\]

\par
